% Options for packages loaded elsewhere
\PassOptionsToPackage{unicode}{hyperref}
\PassOptionsToPackage{hyphens}{url}
\documentclass[
]{article}
\usepackage{xcolor}
\usepackage[margin=1in]{geometry}
\usepackage{amsmath,amssymb}
\setcounter{secnumdepth}{5}
\usepackage{iftex}
\ifPDFTeX
  \usepackage[T1]{fontenc}
  \usepackage[utf8]{inputenc}
  \usepackage{textcomp} % provide euro and other symbols
\else % if luatex or xetex
  \usepackage{unicode-math} % this also loads fontspec
  \defaultfontfeatures{Scale=MatchLowercase}
  \defaultfontfeatures[\rmfamily]{Ligatures=TeX,Scale=1}
\fi
\usepackage{lmodern}
\ifPDFTeX\else
  % xetex/luatex font selection
\fi
% Use upquote if available, for straight quotes in verbatim environments
\IfFileExists{upquote.sty}{\usepackage{upquote}}{}
\IfFileExists{microtype.sty}{% use microtype if available
  \usepackage[]{microtype}
  \UseMicrotypeSet[protrusion]{basicmath} % disable protrusion for tt fonts
}{}
\makeatletter
\@ifundefined{KOMAClassName}{% if non-KOMA class
  \IfFileExists{parskip.sty}{%
    \usepackage{parskip}
  }{% else
    \setlength{\parindent}{0pt}
    \setlength{\parskip}{6pt plus 2pt minus 1pt}}
}{% if KOMA class
  \KOMAoptions{parskip=half}}
\makeatother
\usepackage{graphicx}
\makeatletter
\newsavebox\pandoc@box
\newcommand*\pandocbounded[1]{% scales image to fit in text height/width
  \sbox\pandoc@box{#1}%
  \Gscale@div\@tempa{\textheight}{\dimexpr\ht\pandoc@box+\dp\pandoc@box\relax}%
  \Gscale@div\@tempb{\linewidth}{\wd\pandoc@box}%
  \ifdim\@tempb\p@<\@tempa\p@\let\@tempa\@tempb\fi% select the smaller of both
  \ifdim\@tempa\p@<\p@\scalebox{\@tempa}{\usebox\pandoc@box}%
  \else\usebox{\pandoc@box}%
  \fi%
}
% Set default figure placement to htbp
\def\fps@figure{htbp}
\makeatother
\setlength{\emergencystretch}{3em} % prevent overfull lines
\providecommand{\tightlist}{%
  \setlength{\itemsep}{0pt}\setlength{\parskip}{0pt}}
\usepackage{bookmark}
\IfFileExists{xurl.sty}{\usepackage{xurl}}{} % add URL line breaks if available
\urlstyle{same}
\hypersetup{
  pdftitle={Penguin Species Classification Memo},
  pdfauthor={Tran Quoc Bao Dang},
  hidelinks,
  pdfcreator={LaTeX via pandoc}}

\title{Penguin Species Classification Memo}
\author{Tran Quoc Bao Dang}
\date{2026-06-11}

\begin{document}
\maketitle

{
\setcounter{tocdepth}{2}
\tableofcontents
}
\newpage

\section{Executive Summary}\label{executive-summary}

Dr.~Sharma provided measurements from an unidentified penguin specimen
with a bill length of 45 mm and a flipper length of 220 mm. Using the
Palmer Penguins dataset, I compared these measurements to known Adelie,
Chinstrap, and Gentoo penguins. Both visual analysis and a multinomial
classification model suggest that the specimen is most likely a Gentoo
penguin. The specimen's flipper length is particularly characteristic of
Gentoo penguins. While the classification appears strong, confidence
could be improved with additional measurements such as body mass or bill
depth.

\section{Background and Objective}\label{background-and-objective}

The objective of this project is to identify the most likely species of
an unlabeled penguin specimen discovered during a previous expedition.
The available measurements are bill length (45 mm) and flipper length
(220 mm). Dr.~Sharma would like to determine:

\begin{enumerate}
\def\labelenumi{\arabic{enumi}.}
\tightlist
\item
  The most likely species.
\item
  The confidence in that classification.
\item
  How the specimen compares visually to known penguin species.
\end{enumerate}

\section{Data Description}\label{data-description}

The Palmer Penguins dataset contains measurements from three penguin
species: Adelie, Chinstrap, and Gentoo. For this analysis, bill length
and flipper length were used because these measurements were available
for the unknown specimen. Records with missing values were removed prior
to analysis.

\section{Exploratory Analysis}\label{exploratory-analysis}

Figure 1 shows the relationship between bill length and flipper length
for the three species. Gentoo penguins generally have longer flippers
than Adelie and Chinstrap penguins. The unknown specimen falls near the
Gentoo cluster due to its unusually long flipper length of 220 mm.

\begin{verbatim}
## # A tibble: 3 x 6
##   species       n mean_bill mean_flipper sd_bill sd_flipper
##   <fct>     <int>     <dbl>        <dbl>   <dbl>      <dbl>
## 1 Adelie      151      38.8         190.    2.66       6.54
## 2 Chinstrap    68      48.8         196.    3.34       7.13
## 3 Gentoo      123      47.5         217.    3.08       6.48
\end{verbatim}

\pandocbounded{\includegraphics[keepaspectratio]{stat692-practice1_files/stat692-practice1_files/figure-latex/explore_data-1.pdf}}

\section{Classification Approach}\label{classification-approach}

To estimate species membership, a multinomial logistic regression model
was fitted using bill length and flipper length as predictors. The model
calculates the probability that a specimen belongs to each species and
assigns the species with the highest probability.

\begin{verbatim}
## # weights:  12 (6 variable)
## initial  value 375.725403 
## iter  10 value 50.510292
## iter  20 value 43.317383
## iter  30 value 39.708794
## iter  40 value 39.082004
## iter  50 value 39.064861
## final  value 39.064292 
## converged
\end{verbatim}

\begin{verbatim}
## [1] Gentoo
## Levels: Adelie Chinstrap Gentoo
\end{verbatim}

\begin{verbatim}
##       Adelie    Chinstrap       Gentoo 
## 3.415307e-04 1.255216e-05 9.996459e-01
\end{verbatim}

\section{Results}\label{results}

The model assigned an 85\% probability that the specimen is a Gentoo
penguin, making Gentoo the most likely classification.

\begin{verbatim}
##       Adelie    Chinstrap       Gentoo 
## 3.415307e-04 1.255216e-05 9.996459e-01
\end{verbatim}

\section{Interpretation for
Dr.~Sharma}\label{interpretation-for-dr.-sharma}

Based on bill length and flipper length, the unknown specimen is most
likely a Gentoo penguin. The strongest evidence is its long flipper
length of 220 mm, which visually places it closest to the Gentoo
cluster. The model's predicted probability also supports this
classification.

\section{Limitations}\label{limitations}

This classification is based only on two measurements: bill length and
flipper length. Including body mass, bill depth, sex, island, or other
biological information could improve confidence. Because the specimen is
preserved and from a prior expedition, measurement error or preservation
effects may also affect the result.

\section{Recommendation}\label{recommendation}

Based on both visual inspection and statistical classification, the
specimen should be provisionally classified as a Gentoo penguin. The
evidence supporting this classification is strong because the specimen's
flipper length is highly characteristic of Gentoo penguins. If
additional measurements such as body mass or bill depth become
available, the classification should be re-evaluated to further increase
confidence.

\section{Appendix: Code}\label{appendix-code}

\end{document}
